% §18.5 A Toy Encrypted Search System
%
% Cross-reference resolution status:
%
% DEFINED here:
%   \label{sec:toy-system}                    -- this section
%
% ALREADY DEFINED (chapter 5):
%   \Cref{def:bloom-filter}                   -- chapter 5 §5.2
%
% ALREADY DEFINED (chapter 17):
%   \Cref{def:cipher-map}                     -- §17.1
%
% INTERNAL (chapter 18):
%   \Cref{def:per-document-bloom}             -- §18.2
%   \Cref{sec:confidentiality-measures}       -- §18.3
%   \Cref{sec:adversarial-estimation}         -- §18.4
%   \Cref{sec:closing-the-book}               -- §18.6 (forward)
%
% §18.5 walks through a minimal end-to-end toy encrypted-search
% system, in six items: (1) construction walkthrough (per-document
% Bloom filters, AES-encrypted bit arrays, HMAC tokens, server-side
% bit-vector matching); (2) setup phase; (3) query phase; (4) trace
% through with $\fprate, \fnrate, \pi$ at each stage; (5)
% confidentiality analysis at production-tuned parameters and a
% note about retuning for privacy; (6) notebook hook pointing at
% python/notebooks/ch18-encrypted-search.ipynb.

\section{A Toy Encrypted Search System}
\label{sec:toy-system}

The preceding sections have given the cipher-map framing
(\S\ref{sec:encrypted-search-problem}), the two Bloom-filter
constructions (\S\ref{sec:bloom-based}), the confidentiality
analysis at fixed workload (\S\ref{sec:confidentiality-measures}),
and the estimation problem under adversarial queries
(\S\ref{sec:adversarial-estimation}). This section turns the
framing into running code in the smallest form that exhibits every
piece of the framework: a single-client, single-server,
per-document Bloom system with a standard symmetric cipher for the
encrypted index and a keyed cryptographic function for the query
tokens. The companion notebook materializes the construction; the
prose here gives the structure the notebook fills in.

The system is deliberately minimal. It does not handle multi-user
key management, it does not implement index rotation, it does not
defend against side-channel attacks beyond the cryptographic
primitives' standard guarantees, and it does not optimize for
performance. What it does is exhibit the cipher-map abstraction in
operational form: the client encrypts something, the server
operates on the ciphertext to return an approximate answer, and
the framework prices the resulting trade-off between recovery and
leakage at every stage.

\subsection*{Construction walkthrough}

The system has three cryptographic ingredients. The first is a
keyed cryptographic function $\mathrm{HMAC}_k$ that takes a secret
key $k$ (held only by the client) and a keyword $w$ and produces a
fixed-size query token $\tau_w = \mathrm{HMAC}_k(w)$. The function
is deterministic in $w$ (the same keyword always produces the
same token) and pseudorandom in $w$ from the perspective of any
party that does not hold $k$. The second is a symmetric encryption
scheme (a standard AES-based authenticated encryption, with a key
also held only by the client) that takes a plaintext bit array and
produces a ciphertext blob the server can store but cannot read.
The third is the Bloom filter of \Cref{def:bloom-filter}, used as
the underlying approximate-set structure on each document.

The data structures are equally minimal. Each document $d_i$ is
represented by its keyword set; the client converts the keyword
set into a Bloom filter $B_i$, encrypts the filter as
$\mathsf{Enc}(B_i)$, and uploads $(i, \mathsf{Enc}(B_i))$ to the
server. The server's storage is a dictionary keyed by document ID
and valued by encrypted filter. The client retains the encryption
key and the HMAC key; the server retains the encrypted filter
store; neither party has both the plaintext keywords and the bit
patterns they hash to.

The match operation is the cipher map. To test whether token
$\tau_w$ corresponds to a keyword that was inserted into document
$d_i$, the server reads $\mathsf{Enc}(B_i)$ off disk, decrypts the
filter under the client's key (in the simplest variant; more
sophisticated variants keep the filter encrypted throughout), and
checks whether the bit positions corresponding to $\tau_w$ are all
set in the recovered $B_i$. If yes, the server reports $\top$
(``document $i$ may contain the keyword underlying $\tau_w$''); if
no, the server reports $\bot$ (``document $i$ definitely does not
contain that keyword,'' because Bloom filters have no false
negatives). The simplest variant requires giving the server the
encryption key, which is a security regression; the production
variants use bilinear pairings or one-time-pad-keyed bit arrays to
avoid this, at additional cryptographic cost. For the toy system,
the simplification is acceptable because the goal is to expose the
cipher-map abstraction, not to optimize key management.

\subsection*{Setup phase}

The setup phase runs once per corpus, on the client. The client
holds the plaintext documents $\{d_1, \ldots, d_n\}$, the
encryption key $K_{\mathrm{enc}}$, and the HMAC key $K_{\mathrm{tok}}$.
For each document, the client extracts the keyword set, computes a
Bloom filter at the target $\fprate$ (the chapter~5 formula
$m \approx 1.44 \cdot W_d \cdot \log_2(1/\fprate)$ gives the
required filter size), encrypts the filter, and uploads the
encrypted blob to the server with the document's identifier. The
upload is a single round-trip per document; the total bandwidth is
$n \cdot m / 8$ bytes (filter size times document count).

After setup, the server holds the encrypted index but cannot read
any plaintext keyword, and the client holds the two keys but no
copy of the encrypted index. The corpus is in deployment.

\subsection*{Query phase}

A query is a three-step protocol. The client computes the query
token $\tau_w = \mathrm{HMAC}_{K_{\mathrm{tok}}}(w)$ locally and
sends $\tau_w$ to the server. The server runs the match operation
against every encrypted filter in the store and assembles the set
of document IDs for which the match succeeded. The server returns
the matched ID set to the client; the client verifies the
candidates (typically by a second-pass decryption of the matched
documents' contents, which is outside the encrypted-search layer)
and discards the false-positive matches. The protocol completes in
two round-trips: one for the token, one for the matched IDs.

The server never sees the plaintext keyword $w$ (the HMAC hides
it), never sees the plaintext keyword set of any document (the
Bloom filter hides them up to false positives), and never sees the
client's verification step (it is entirely client-side). What the
server sees is a sequence of query tokens, a sequence of match
decisions, and the bare fact of which documents matched which
tokens. The information flow from the query phase is exactly the
cipher-map output stream the chapter~17 framework predicted, and
the chapter~7 PPV analysis prices it.

\subsection*{Trace through a single query}

Take a query for the keyword \texttt{alice} against a corpus of
$|D| = 1000$ documents, with per-document Bloom filters at
$\fprate = 0.01$ and the framework's $\fnrate = 0$. Suppose the
true match set is $\{17, 209, 832\}$, three documents containing
the queried keyword. The trace proceeds as follows.

\begin{enumerate}
\item Client side. The client computes $\tau_{\mathrm{alice}} =
  \mathrm{HMAC}_{K_{\mathrm{tok}}}(\mathtt{alice})$, a fixed-size
  pseudorandom bit string. The framework's $\fprate$ and
  $\fnrate$ enter through the filter parameters built into the
  encrypted filters during setup; they do not affect the token
  computation itself. The framework's $\pi$ is fixed by the
  corpus and the query: in this case $\pi = 3/1000 = 0.003$,
  three of the thousand documents truly contain \texttt{alice}.

\item Server side. The server tests $\tau_{\mathrm{alice}}$
  against each of the 1000 encrypted filters. The three true
  matches return $\top$ unconditionally (the Bloom filter
  $\fnrate = 0$ guarantees this). The remaining 997 non-matches
  return $\top$ with probability $\fprate = 0.01$; the expected
  number of false-positive matches is $997 \cdot 0.01 \approx
  9.97$. The server's reported match set has expected size
  $3 + 9.97 \approx 13$ documents.

\item Return path. The server returns the matched ID set (about
  13 IDs on average) to the client. The client decrypts the
  matched documents (or queries a higher-tier verification
  service), confirms the three true matches, and discards the
  remaining ten as false positives. The framework's $\PPV$ for
  this query is
  \[
    \PPV
    \;=\; \frac{0.003}{0.003 + 0.997 \cdot 0.01}
    \;\approx\; \frac{0.003}{0.01297}
    \;\approx\; 0.231,
  \]
  so $C = 1 - \PPV \approx 0.769$. The system's matches at this
  $\pi$ and this $\fprate$ are dominated by false positives,
  and the confidentiality margin is correspondingly large.
\end{enumerate}

The trace exhibits every framework parameter. The $\fprate$ is
the per-document false-positive rate of the Bloom filter; the
$\fnrate = 0$ is the structural property the construction
inherits; the $\pi$ is the workload prevalence of the queried
keyword; the resulting $\PPV$ and $C$ price the trade-off. The
framework predicted these quantities before the trace ran, and
the trace exhibits them in operational form.

\subsection*{Confidentiality analysis at production-tuned
parameters}

The toy system as built is tuned for correctness, not for
privacy. A more representative production trace uses a higher
workload prevalence (the encrypted-search literature typically
analyzes deployments where the queried keywords are
\emph{interesting} keywords, not arbitrary vocabulary items).
At $\fprate = 0.01$ and $\pi = 0.1$ (one queried keyword in ten
truly matches the queried document) with 10 keywords per
document, the system has
\[
  \PPV
  \;=\; \frac{0.1}{0.1 + 0.9 \cdot 0.01}
  \;\approx\; 0.917,
\]
which puts $C \approx 0.083$. A server compromise at this
operating point exposes most of the plaintext: an adversary
reading match decisions infers the underlying keyword with
about a $92\%$ posterior probability. The toy system is not
privacy-tuned, and a deployment with serious privacy
requirements has to retune $\fprate$ upward to push the
operating point further onto the confidentiality side of the
trade-off.

The retuning is mechanical. Solving the chapter~7 PPV equation
for $\fprate^*$ at the target $C$ gives $\fprate^* = \pi C /
[(1 - \pi)(1 - C)]$. For a target $C = 0.5$ at $\pi = 0.1$,
the required $\fprate^*$ is
\[
  \fprate^*
  \;=\; \frac{0.1 \cdot 0.5}{0.9 \cdot 0.5}
  \;=\; \frac{0.05}{0.45}
  \;\approx\; 0.111.
\]
The system has to allow about an $11\%$ per-document false
positive rate to achieve fifty-fifty leakage. The filter size
shrinks correspondingly (chapter~5's
$m \approx 1.44 \cdot W_d \cdot \log_2(1/\fprate)$ formula
gives $m \approx 1.44 \cdot 10 \cdot 3.17 \approx 46$ bits per
filter, against the $\fprate = 0.01$ value of $m \approx 96$
bits), and the client's post-filtering load grows
correspondingly. The framework's space-accuracy-confidentiality
triangle (chapter~8 set the entropy floor, chapter~17 priced
the confidentiality) tells the designer that no choice escapes
the trade-off; the only design freedom is which corner of the
triangle to live in.

\subsection*{Notebook hook}

The companion notebook for this chapter,
\texttt{python/notebooks/ch18-encrypted-search.ipynb}, implements
the toy system end-to-end. The notebook builds the per-document
Bloom filters, encrypts them with a standard library AES
primitive, computes HMAC-based query tokens, runs the server-side
match, and verifies the recovered match sets against the ground
truth. Plots compare the empirical $\PPV$ and confidentiality
across $\fprate$ values from $0.001$ to $0.5$, exhibiting the
chapter's trade-off curve as a measured quantity rather than a
theoretical prediction. The notebook also runs the
\Cref{sec:adversarial-estimation} comparison between iid Chernoff
and MAB UCB confidence intervals on a synthetic adversarial query
trace, exhibiting the two-orders-of-magnitude gap the worked
example predicted.

The notebook is intended as a starting point, not as a production
artifact. Readers interested in deploying an encrypted-search
system should consult the production-grade implementations in
the SSE literature and the underlying cryptographic primitives'
formal security definitions; the notebook here exhibits the
Bernoulli framework on the construction without engaging the
full cryptographic-engineering surface.
